\subsection{Search and Minimum-Propellant Formulation}
We used a feedback controller on the planar model to generate starting trajectories. It varied the ignition altitude, initial inventory, pitch bias, and descent-speed profile. During the burn, the pitch target changes from a biased retrograde direction to upright with lateral-velocity feedback. The throttle tracks a decreasing descent speed, and the gimbal supplies a pitch-tracking moment.

The optimization minimizes $m(0)-m_d$ over phase durations, throttle and gimbal commands, and trajectory states. All seven states are continuous across phase boundaries. Coast phases have zero thrust. Powered phases obey $0.4\leq u_T\leq1$ and $|\delta|\leq10^\circ$. The ship position is free during optimization and fixed during replay. The terminal targets are deck contact, $v_x=0$, $v_h=-1.5\,\mathrm{m/s}$, upright attitude, and zero pitch rate. The optimizer includes 75 kg above the 5 t usable reserve to allow for integration differences in the replay.

Hermite--Simpson direct collocation initializes the planar solutions~\cite{kelly}. RK4 multiple shooting refines them and reoptimizes the spherical design. Throttle and gimbal are linearly interpolated within each phase. Each shooting interval uses eight RK4 substeps, with a local trust region on phase duration and pitch. CasADi supplies derivatives to IPOPT~\cite{casadi,ipopt}. We then integrate each accepted control schedule independently to first contact and check the touchdown limits at the selected ship position.

\subsection{Planar Comparison and Inventory Accounting}
At $15^\circ$ the planar model gives 26.23 t and a ship location of 278.31 km with one ignition. The lowest verified planar two-ignition result requires 27.61 t and drives its first pulse to the 0.001 s numerical duration floor. A feedback trajectory with a separate 1.5 s entry burn lands with 38.65 t. The one-ignition planar result seeds the spherical reference in Section~\ref{sec:spherical-results}.

We account for the separation inventory as
\begin{equation}\label{eq:p2-propellant-budget}
 m_{p,0}=m_{\mathrm{burned}}+m_{\mathrm{reserve}}+m_{\mathrm{residual}}+m_{\mathrm{surplus}}.
\end{equation}
The nominal spherical solution burns 20.89 t and retains 5.08 t, of which 5 t is the planned usable reserve and about 0.08 t is surplus. The unusable residual is set to zero.
